\section{Introduction to API Decorators}

\slidebackground{backgrounds/wallpaper.jpg}

\begin{frame}{What Are API Decorators?}
	\begin{itemize}
		\item In Odoo, methods are often decorated with \textbf{\texttt{@api....}} markers.
		\item A decorator tells Odoo \textbf{how} the method should be called and when it should run.
		\item Decorators are not just Python sugar -- they change ORM behavior.
	\end{itemize}
\end{frame}

\begin{frame}{Common Decorators We Will Cover}
	\begin{columns}[T]
		\begin{column}{0.5\textwidth}
			\begin{itemize}
				\item \texttt{@api.model}
				\item \texttt{@api.model\_create\_multi}
				\item \texttt{@api.depends}
				\item \texttt{@api.onchange}
			\end{itemize}
		\end{column}
		\begin{column}{0.5\textwidth}
			\begin{itemize}
				\item \texttt{@api.constrains}
				\item \texttt{@api.ondelete}
				\item \texttt{@api.autovacuum}
				\item \texttt{@api.private} / \texttt{@api.readonly}
			\end{itemize}
		\end{column}
	\end{columns}
\end{frame}

\begin{frame}{Why Decorators Matter}
	\begin{itemize}
		\item They decide whether a method is model-level or record-level.
		\item They connect computed fields to their dependencies.
		\item They validate data before it is saved.
		\item They control UI onchange behavior in forms.
		\item Using the wrong decorator creates subtle bugs.
	\end{itemize}
\end{frame}
